\subsection{Eliminating every nonsimplicial rank-eight endpoint cell}
\label{subsec:k12-r8-nonsimplicial}

The preceding excess-five reduction can be completed in rank eight,
including the possible ambient glue. The additional inputs are the
classification of short perfection relations and the minimum-basis
index bounds. These are invoked as cited theorems, with the later
corrections and completion of the index bound included.

\begin{theorem}[K17-K12-R8-NONSIMPLICIAL]
\label{thm:k12-r8-nonsimplicial}
Let $L$ be a well-rounded lattice of rank eight. Suppose the cone
$\sigma$ generated by all its minimum quadratic rays has dimension
thirteen and is nonsimplicial. Then $\sigma$ has an integral
orientation-reversing stabilizer.

Consequently every nonsimplicial rank-eight projective degree-twelve
Voronoi cell is zero in oriented coinvariants over any ring in which
two is invertible. There is no assumption that the minima generate
the ambient lattice over $\mathbb Z$.
\end{theorem}

\begin{proof}
Normalize the minimum to one. We first reconstruct the equality case
of the circuit argument to specify the core to which the source
classification will apply. Choose a circuit of $k$ distinct minimum
quadratic rays and let $W$ be the span of their underlying vectors,
of dimension $m$. Its quadratic rank is $k-1$. Splitting its relation
into positive and negative parts gives equal positive semidefinite
sums. Equality of their kernels shows that each sign support spans
$W$, so $k\geq2m$.

Let $\sigma_W$ be the cone of all minimum rays in $W$, of dimension
$r_W$. It is the actual perfect cone for $L\cap W$, and it is a face
of $\sigma$, exposed by a positive semidefinite functional with
kernel $W$. It is nonsimplicial: distinct rank-one positive
semidefinite rays are extremal in a cone generated by such rays,
and the chosen circuit lies among them. Choose $8-m$ minima
independent modulo $W$. Projecting their quadratic rays gives
\[
 13\geq r_W+8-m,
 \qquad r_W\geq k-1\geq2m-1.
\]
Thus $m\leq6$ and $r_W\leq m+5\leq11$.
By \cite[Theorem~1.1(ii)]{DutourHulekSchurmann2015}, the only
nonsimplicial perfect cone of dimension at most ten is the $D_4$
cone. Its vector rank four and quadratic rank ten contradict
$r_W\leq m+5$. Hence
\[
             m=6,\qquad r_W=11,\qquad k=12.             \tag{R8.1}
\]
Each sign support therefore consists of six independent minima.

Choose minima $v,w$ independent modulo $W$. The classes of
$vv^t,ww^t$ form a basis of the quotient of the thirteen-dimensional
quadratic span by the eleven-dimensional core span. Split the vector
space as $W\oplus U$, using the images of $v,w$ as a basis of $U$.
For another outside minimum $x$, with $U$-coordinates $(a,b)$,
comparison of its $U\times U$ quadratic block with a combination
of $vv^t,ww^t$ forces $ab=0$. If $a\ne0$, the $W\times U$ block
then gives $x=av$; the other case gives $x=bw$. Equality of the
minimum norms makes the nonzero scalar $\pm1$. Thus $v,w$ are the
only two outside minimum pairs and are rational coloops.

\medskip\noindent
\emph{The integral six-dimensional core.}
Let $M_0$ be generated by the twelve circuit vectors. It has minimum
one, the two sign supports are independent six-element systems of
minima, and their union generates $M_0$. The circuit is
perf-irreducible in the sense of
\cite[Definition~2.4 and \S3]{BergeMartinet2009}: a splitting into
two perfection relations would contradict its minimality.
The hypotheses of \cite[Corollary~6.9]{BergeMartinet2009} therefore
hold. In rank six that result leaves $D_6$ or the Watson
configuration of \cite[Theorem~5.1]{BergeMartinet2009}.

The $D_6$ alternative is impossible. Its thirty minimum pairs would
also be minimum pairs of $L$, and their quadratic rays span
$\operatorname{Sym}^2(W)$ of dimension twenty-one, contradicting
(R8.1). Indeed, the roots $e_i+e_j,e_i-e_j$ span the off-diagonal
quadratic coordinates by subtraction and the diagonal coordinates
by their sums. We may therefore choose independent minima
$f_1,\ldots,f_6$ such that
\[
 e=\frac{f_1+\cdots+f_6}{3},\qquad
 M_0=\sum_{i=1}^6\mathbb Zf_i+\mathbb Ze,
 \qquad
 \text{circuit pairs }\{f_i,e-f_i:1\leq i\leq6\}.        \tag{R8.2}
\]
Set $F_0=\sum_i\mathbb Zf_i$ and $M=L\cap W$. The frame has
$[M_0:F_0]=3$. The six-dimensional minimum-basis index bound in
\cite[Theorem~5.2]{MartinetIndex2001} gives
\[
                   3[M:M_0]=[M:F_0]\leq4.
\]
Hence $M=M_0$. Saturation of the circuit lattice has been proved,
rather than assumed in the application of the short-relation theorem.

\medskip\noindent
\emph{The full minimum set of the core.}
Use the $f_i$ as rational coordinates and let $S$ be the span of
$f_if_i^t$ and $(e-f_i)(e-f_i)^t$. The circuit gives $\dim S=11$,
so this is the entire span of $\sigma_W$. If a nonzero vector $x$
has $xx^t\in S$, its off-diagonal coordinates satisfy
\[
 x_jx_k=c-d_j-d_k\quad(j\ne k),
 \qquad c=\frac13\sum_{i=1}^6d_i.                       \tag{R8.3}
\]
In fact, if the coefficient of $(e-f_i)^2$ is $b_i$, take
$d_i=b_i/3$; the other six rays contribute only to the diagonal.
Alternating four entries of (R8.3), for distinct $j,k,l,m$, gives
\[
                       (x_j-x_m)(x_k-x_l)=0.
\]
If $x_a\ne x_b$, every coordinate outside $\{a,b\}$ has a common
value. If both $x_a,x_b$ differed from this value, applying the last
identity with two outside indices would contradict it. Thus all but
at most one coordinate of $x$ coincide. Write $x=a\mathbf1+bf_i$.
Equation (R8.3) forces $d_j=d$ for $j\ne i$ and $d_i=d-ab$, hence
$c=2d-ab/3$ and
\[
                         a^2=-ab/3.
\]
If $a=0$, $x$ lies on the line of $f_i$. Otherwise $b=-3a$,
so it lies on the line of $e-f_i$. The all-equal case gives only
zero. Consequently the twelve lines in (R8.2) are the only nonzero
rank-one lines in $S$. Equal minimum norms show that their signed
representatives are the complete core minimum set.

\medskip\noindent
\emph{Core symmetries and orientation.}
Every permutation $\pi$ of the $f_i$ fixes $e$ and preserves $M$
and its twelve minimum pairs. On the twelve-dimensional free ray
space it acts by two copies of $\pi$, with determinant
$\operatorname{sgn}(\pi)^2=1$. It fixes the unique relation
\[
             \sum_i f_if_i^t-\sum_i(e-f_i)(e-f_i)^t=0.   \tag{R8.4}
\]
Therefore its determinant on the eleven-dimensional quotient $S$
is $+1$.

Define $F(f_i)=e-f_i$. Then $F(e)=e$ and $F^2=1$, so $F$ is an
integral automorphism of $M$. It exchanges the two ray systems by
six transpositions, giving determinant $+1$ on the free ray space,
and negates (R8.4). Its determinant on $S$ is consequently $-1$.
Every $\pi F$ also reverses orientation. Extending such a map by
the identity on $v,w$ keeps this sign. Swapping $v,w$ while fixing
$M$ has sign $-1$ whenever it preserves the ambient lattice.
These are cone automorphisms; no orthogonality of the original
witness form is needed.

\medskip\noindent
\emph{Bounding the ambient glue.}
The lattice generated by all minima is
\[
                  L_{\min}=M\oplus\mathbb Zv\oplus\mathbb Zw.
\]
Write $g=[L:L_{\min}]$. The independent minimum frame
$f_1,\ldots,f_6,v,w$ has index $3g$ in $L$.
The completed rank-eight bound in
\cite[appendix, pp.~6--8]{MartinetIndexErrata2014} gives index at
most nine outside $E_8$, and
\cite[Theorem~10.1]{MartinetIndex2001} shows that index nine itself
occurs only for $E_8$. The latter lattice is impossible here: its
complete minimum quadratic span has dimension thirty-six, whereas
$\dim\sigma=13$. It follows that
\[
                           3g\leq8,
                  \qquad g\in\{1,2\}.                  \tag{R8.5}
\]
If $g=1$, extend $F$ by fixing $v,w$ to obtain the desired odd
integral stabilizer.

\medskip\noindent
\emph{Every admissible binary glue.}
For $g=2$, write
\[
 L=L_{\min}+\mathbb Z\frac{u+av+bw}{2},
       \qquad u\in M,\quad a,b\in\{0,1\}.
\]
The outside bits cannot both vanish, since $L\cap W=M$.
If $a=b=1$, interchanging $v,w$ and fixing $M$ preserves the
glue class and reverses orientation.

Suppose exactly one outside bit is one, say $a=1,b=0$.
Since $[M:F_0]=3$ is odd, the six $f_i$ form a basis of $M/2M$.
Represent $u$ by a binary word of weight $k$. Modulo $2M$ one has
\[
 e=\sum_i f_i,
 \qquad F(u)=u+k\sum_i f_i.
\]
For even $k$, $F$ fixes the glue class. For $k=3$, it complements
the six bits; choose a permutation $\pi$ sending the complementary
three-element support back to the original support. Then $\pi F$
fixes the class. In both cases the extension by the identity on the
outside directions is an integral odd stabilizer of $L$.

The remaining weights $k=1,5$ are impossible. A weight-one class
is represented by a minimum $f_i$, and a weight-five class by a
minimum $e-f_i$. Hence $u\equiv m\pmod{2M}$ for some core minimum
$m$. Both $(m+v)/2$ and $(m-v)/2$ belong to $L$ and are nonzero,
while
\[
 Q\!\left[\frac{m+v}{2}\right]
 +Q\!\left[\frac{m-v}{2}\right]
       =\frac{Q[m]+Q[v]}2=1.
\]
This contradicts their having norms at least one. The other
one-bit case is identical. All possible ambient glues have now
been treated, proving the theorem.
\end{proof}

\begin{corollary}
\label{cor:k12-r8-simplicial-gate}
In the localized rank-eight projective degree-twelve Voronoi
chain group, only simplicial cells can contribute. Equivalently,
each surviving cell in that degree has exactly thirteen minimum
quadratic rays.
\end{corollary}

\paragraph{Certificate and retained inputs.}
The complete standard-library verifier
\texttt{k12\_rank8\_watson\_certificate.py} checks the integral
Watson model, its frame index three and quadratic rank eleven,
the quotient orientation sign, and all $192$ binary glue classes
with nonzero outside bits. It supplies odd maps for $168$ classes
and minimum representatives proving the other $24$ impossible.
Its JSON output is retained. The no-extra-rank-one argument above
is symbolic and does not depend on finite sampling.

The short-relation and minimum-index classifications remain
explicit cited inputs. The known EVGS rank-six data provide a
separate corroboration: among $329$ codimension-ten records there
is exactly one with twelve pairs, the Watson model. This database
check is not required by the proof and does not rerun its coverage.
The remaining simplicial rank-eight cells and their boundary
matrices have not been eliminated. In particular this theorem
does not assert the rank-eight $H_5$ vanishing or
$K_{12}(\mathbb Z)=0$.
